% © 2026 Ing. David Jaroš — CC BY-NC-ND 4.0
%
% File: research_tracks/BRIDGE_CLOSURE/G137B_AUDIT/A1_complete_spin_structure_trace_theorem.tex
% Status: A1 bridge audit — complete spin-structure trace theorem.
% Purpose: Defend why the compact Layer2 measure uses eta^{-2} theta2 theta3 theta4.

\documentclass[12pt,a4paper]{article}
\usepackage[a4paper,margin=2.5cm]{geometry}
\usepackage{amsmath,amssymb,amsthm,mathtools,booktabs,longtable}
\usepackage[most]{tcolorbox}
\usepackage{hyperref}
\usepackage{xcolor}
\usepackage{microtype}

\hypersetup{colorlinks=true,linkcolor=blue!60!black,urlcolor=blue!60!black,citecolor=green!50!black}

\newtheorem{definition}{Definition}[section]
\newtheorem{lemma}[definition]{Lemma}
\newtheorem{proposition}[definition]{Proposition}
\newtheorem{theorem}[definition]{Theorem}
\newtheorem{remark}[definition]{Remark}

\newtcolorbox{statusbox}[1][]{
  colback=yellow!8!white,
  colframe=orange!70!black,
  arc=3pt,
  boxrule=1pt,
  title=Audit Status,
  #1
}

\title{A1 Audit: Complete Spin-Structure Trace Theorem for G137-B}
\author{Ing. David Jaroš}
\date{2026-06-14}

\begin{document}
\maketitle

\begin{abstract}
This note audits the most sensitive bridge in the final G137-B closure: why the
compact one-real Layer2 determinant should contain the complete product
\[
  \eta^{-2}\vartheta_2\vartheta_3\vartheta_4
\]
rather than a single spin sector, a sum of sectors, or a hand-picked modular
factor.  The theorem proved here is conditional on the current UBT compact-sector
bridges: SU(2) Scherk--Schwarz splitting, odd-winding parity measure, and
quaternionic triad completeness of $\mathrm{Im}\,\mathbb H$.  Under these
conditions, the minimal multiplicative readout trace closed under the three
nontrivial torus spin/readout structures is uniquely
$\vartheta_2\vartheta_3\vartheta_4$, with two real oscillator denominators
$\eta^{-2}$.  Therefore
\[
  Z_{\rm compact}^{1\rm real}(\tau)
  =\eta^{-2}\vartheta_2\vartheta_3\vartheta_4=2\eta(\tau)
\]
by Jacobi's identity.  This upgrades the trace from an ad hoc insertion to a
minimal complete compact-sector measure.
\end{abstract}

\begin{statusbox}
\textbf{A1 result.} Complete spin-structure trace is defended as the minimal
multiplicative trace compatible with SU(2) twist + odd parity + quaternionic
triad completeness.\
\textbf{Remaining attack surface.} A reviewer can still challenge whether
triad completeness necessarily maps to the three torus theta structures.  That
is now the precise bridge to audit.
\end{statusbox}

\section{Bridge under audit}

The final G137-B closure uses
\[
  Z_{\rm compact}^{1\rm real}(\tau)
  =\eta(\tau)^{-2}\vartheta_2(0|\tau)\vartheta_3(0|\tau)\vartheta_4(0|\tau).
\]
The audit question is:
\begin{quote}
Why must the compact Layer2 trace contain all three theta blocks
$\vartheta_2$, $\vartheta_3$, and $\vartheta_4$ multiplicatively?
\end{quote}

\section{Three sources of the three theta blocks}

The current UBT compact-sector bridges provide three distinct sectors:

\begin{longtable}{llll}
\toprule
UBT source & Compact/readout meaning & Theta block & Status \\
\midrule
\endhead
Neutral diagonal fields $a,d$ & periodic integer winding & $\vartheta_3$ & SU(2) twist bridge \\
Charged off-diagonal fields $b,c$ & shifted / anti-periodic tower & $\vartheta_2$ & SU(2) Scherk--Schwarz bridge \\
Odd-winding parity trace & $(-1)^n$ weighted integer tower & $\vartheta_4$ & odd-winding parity bridge \\
\bottomrule
\end{longtable}

Thus each theta block has an independent structural origin.  None is introduced
only to repair the coefficient.

\section{Why product, not sum}

For Gaussian compact sectors, independent determinant factors multiply.  If the
compact readout block decomposes into independent periodic, shifted, and parity
subsectors, then the one-real determinant is of the form
\[
  Z_{\rm compact}=Z_{\rm osc}^{-1}Z_{\vartheta_2}Z_{\vartheta_3}Z_{\vartheta_4}.
\]
A sum would correspond to an incoherent statistical mixture of mutually
exclusive sectors.  The UBT Layer2 readout instead uses simultaneous compact
constraints from the full $\mathrm{Im}\,\mathbb H$ triad.  Therefore the trace is
multiplicative.

\begin{lemma}[Multiplicativity of independent Gaussian blocks]
If the compact quadratic operator factorises as
\[
  P=P_2\oplus P_3\oplus P_4,
\]
then its determinant factorises as
\[
  \det P=(\det P_2)(\det P_3)(\det P_4),
\]
and the partition function is multiplicative.
\end{lemma}

\begin{proof}
For block-diagonal operators, the spectrum is the union of the spectra of the
blocks.  The zeta-regularised determinant is exponential in the spectral zeta
sum, hence additive in logarithms and multiplicative in determinants.
\end{proof}

\section{Why complete trace, not one sector}

The imaginary-quaternion sector is a triad:
\[
  \mathrm{Im}\,\mathbb H=\mathrm{span}\{I,J,K\}.
\]
A compact Layer2 trace that keeps one of the three torus spin/readout structures
but discards the others would select a preferred element of this triad.  That
breaks the internal quaternionic symmetry of the compact phase sector.

\begin{definition}[Triad-complete compact trace]
A compact Layer2 trace is triad-complete if it includes one representative of
each of the three nontrivial readout sectors associated with the
$\mathrm{Im}\,\mathbb H$ phase triad.
\end{definition}

\begin{theorem}[Minimal complete spin-structure trace]
Assume:
\begin{enumerate}
\item SU(2) Scherk--Schwarz splitting gives periodic and shifted compact towers;
\item odd-winding parity supplies the $(-1)^n$ weighted trace;
\item Layer2 compact readout is triad-complete under $\mathrm{Im}\,\mathbb H$;
\item independent Gaussian compact blocks multiply.
\end{enumerate}
Then the minimal compact one-real Layer2 trace is
\[
  Z_{\rm compact}^{1\rm real}(\tau)
  =\eta(\tau)^{-2}\vartheta_2(0|\tau)\vartheta_3(0|\tau)\vartheta_4(0|\tau).
\]
\end{theorem}

\begin{proof}
The SU(2) twist supplies $\vartheta_3$ and $\vartheta_2$.  Odd-winding parity
supplies $\vartheta_4$.  Triad completeness requires all three sectors; minimality
requires each once.  Gaussian independence makes the determinant multiplicative.
The two compact real oscillator directions contribute $\eta^{-2}$.  Therefore
the stated product is the minimal complete trace.
\end{proof}

\section{Immediate coefficient consequence}

Using Jacobi's identity,
\[
  \vartheta_2\vartheta_3\vartheta_4=2\eta^3,
\]
we obtain
\[
  Z_{\rm compact}^{1\rm real}(\tau)
  =\eta^{-2}\vartheta_2\vartheta_3\vartheta_4
  =2\eta(\tau).
\]
Thus the coefficient factor in G137-Bb follows:
\[
  B(\rho)=12^{3/2}(2\eta(i\rho))^{1/4}.
\]

\section{Reviewer attack table}

\begin{longtable}{p{0.28\textwidth}p{0.32\textwidth}p{0.30\textwidth}}
\toprule
Attack & Response & Residual risk \\
\midrule
\endhead
Why not only $\vartheta_2/\eta$? & That keeps only the shifted charged sector and drops periodic and parity traces. & Would be valid only if Layer2 trace were not triad-complete. \\
Why not orbifold sum? & A sum is a mixture of alternatives; independent Gaussian constraints multiply. & Requires accepting determinant factorisation. \\
Why include $\vartheta_4$? & Odd-winding parity bridge supplies the $(-1)^n$ trace. & Depends on odd-winding bridge audit. \\
Why include all three once? & Minimal triad-complete trace under $\mathrm{Im}\,\mathbb H$. & Triad-to-spin-structure map is the key external audit point. \\
\bottomrule
\end{longtable}

\section{Conclusion}

A1 is not a new numerical fit.  It is an audit theorem: the complete product
$\vartheta_2\vartheta_3\vartheta_4$ is the minimal multiplicative compact trace
compatible with the current UBT compact-sector bridges.  The next bridge audits
should attack the two dependencies that remain most exposed: the four-channel
geometric mean and the SU(2) twist uniqueness.

\end{document}
